\input{\string~/Documents/School/"UC Irvine"/ta_template2.0.tex}
\rh{2B}{5.1}
\chead{\today}
\graphicspath{ {\string~/Desktop} }
\usetikzlibrary{arrows}
%============ANSWER KEY TOGGLE===========
\newcommand{\ans}[1]{\noindent {\color{red} \textbf{Answer:} #1}}
\renewcommand{\ans}[1]{\vspace{3in}}
%==========================================

\begin{document}

\begin{center}
\textbf{Areas \& Distances}
\end{center}

The main idea of this section is that we can try to approximate areas of non-standard regions of the plane using rectangles, and then adding more rectangles to decrease error. First, we'll review sigma notation:

\begin{aun}[Sigma Notation]
Sums can be concisely expressed using notation in which one utilizes an \emph{index}, or variable that takes on consecutive number values. For example:
\[
\sum_{i = 4}^7 2i = 2(4) + 2(5) + 2(6) + 2(7) = 8 + 10 + 12 + 14 = 44
\]
Recall that distributivity (pulling out constants) and associativity (splitting sums into two sums) both work. Also:
\[
\sum_{j= 1}^n j = \frac{n (n+1)}{2} \mte{,} \sum_{j = 1} ^ {n} j^2 = \frac{n (n+1)(2n+1)}{6} \mte{ and } \sum_{j= 1}^n j^3 = \lrp{\frac{n (n+1)}{2}}^2
\]
\end{aun}

\begin{defi}[Riemann Sum]
Let $f(x)$ be a continuous function. we \emph{approximate} the area between $f(x)$ and the $x$-axis on the interval $[a,b]$ using $n$ rectangles via:
\[
\sum_{i = 1}^n f(x_i^*) \Delta x 
\]
where:
\[
x_i^* = \begin{cases}
a + (i-1) \Delta x & \text{ for Left-hand Sum, written } L_n\\
a + i \Delta x & \text{ for Right-hand Sum, written } R_n\\
a + \frac{i}{2} \Delta x & \text{ for Midpoint Sum, written } M_n\\
\end{cases}
\mte{ and } \Delta x = \frac{b - a}{n}
\]
\end{defi}
\begin{defi}
To compute the \emph{exact area} of this region between $f(x)$ and the $x$-axis on the interval $[a,b]$, we use:
\[
A = \lim_{n \to \infty} L_n =  \lim_{n \to \infty} R_n = 
\]
\end{defi}


\begin{exx}
Expand the following sums:
	\[
	\sum_{i=1}^7 i  \mte{ and } \sum_{j=4}^7 (j^2 +k)
	\]
\end{exx}
\ans{\emph{Solution:} 
	\[
	\sum_{i=1}^7 i = 1 + 2 + \cdots + 7 = 28
	\]
	and
	\[
	\sum_{j=4}^7 (j^2 +k) = 16 + \cdots + 49 + 4k
	\]}
\problembox{Complete the following parts of the problem with a friend:
		\begin{enumerate}
		\item Sketch the graph of $f(x) = \frac{1}{x}$ from $a = 1$ to $b = 2$ 
		\item Compute the integral $\int \frac{1}{x}dx$ with four approximating rectangles using right end points. {\color{gray} HINT: Use the picture from above!}
		\item State whether your answer is an underestimation or an overestimation.
		\end{enumerate}}\\
\ans{\href{https://www.desmos.com/calculator/ynmwlxf1cv}{(Graph)}}

\problembox{Set up but \textbf{do not evaluate} a limit to approximate the area of the region between the function and the $x$-axis of:
\[
f(x) := \frac{2x}{x^2+1} \mte{ on the interval } [1,3]
\]
}

\ans{We can take a right hand limit, so let's set up $R_n$ first:
\[
\Delta x = \frac{3-1}{n} = \frac2n \mte{ and } x_i^* = 1 + \frac{2i}{n}
\]
and just plug in:
\[
R_n = \sum_{i = 1}^n f\lrp{1 + \frac{2i}{n}}\cdot \frac{2}{n} = \sum_{i = 1}^n \frac{2 \lrp{\frac{2i}{n}}}{\lrp{\frac{2i}{n}}^2 + 1} \cdot \frac{2}{n}
\]
To complete our answer, just add a limit in:
\[
A = \lim_{n \to \infty} R_n =  \lim_{n \to \infty} \sum_{i = 1}^n \frac{2 \lrp{1 + \frac{2i}{n}}}{\lrp{1 + \frac{2i}{n}}^2 + 1} \cdot \frac{2}{n}
\]
Phew.
}

\end{document}